% ============================================================
%  CHAPTER 5 — Experiments and Results
%  Logical chain: Validation → Results → Comparison
%  Status: Not started | Metrics placeholders (PQ6)
% ============================================================

\chapter{Experiments and Results}
\label{chap:experiments}

% Opening: restate evaluation goals; connect to H1, H2, H3.
% Remind the reader that all experiments are simulation-only (scope, §1.5).

\lipsum[31]

% ─── 5.1 EXPERIMENTAL SETUP ─────────────────────────────────
\section{Experimental Setup}
\label{sec:exp-setup}

% Summarize simulator, robot model, control parameters.
% Reference Ch. 4 — do not repeat details, only add what is needed for
% experiment reproducibility.

\lipsum[32]

% ─── 5.2 METRICS ────────────────────────────────────────────
\section{Evaluation Metrics}
\label{sec:metrics}

% All thresholds (X, Y, Z) to be grounded in cited literature (PQ6).
% Until resolved, keep as parameters.

\begin{table}[htb]
  \centering
  \caption{Evaluation metrics (thresholds to be grounded in literature, PQ6).}
  \label{tab:metrics}
  \begin{tabularx}{\linewidth}{llXl}
    \toprule
    \textbf{Metric} & \textbf{Symbol} & \textbf{Definition} & \textbf{Target} \\
    \midrule
    CoM tracking error       & $e_{\mathrm{CoM}}$   & RMSE of $\com$ vs.\ reference & $< X$~cm \\
    Contact force compliance  & $\eta_{\mathrm{FC}}$ & Fraction of time inside friction cone & $> 95\%$ \\
    Task success rate         & $\eta_{\mathrm{task}}$ & \% of trials completing target task & $> Y\%$ \\
    QP solve time             & $t_{\mathrm{QP}}$    & Mean/max per control cycle & $<1,kHz$ \\
    Perturbation recovery     & $\eta_{\mathrm{rec}}$ & Recovery rate after impulse & TBD \\
    End-effector error        & $e_{\mathrm{EE}}$    & Position RMSE during manipulation & $< Z$~cm \\
    \bottomrule
  \end{tabularx}
\end{table}

% ─── 5.3 BASELINES ──────────────────────────────────────────
\section{Baselines}
\label{sec:baselines}

\begin{enumerate}
  \item \textbf{Null-space Task-Priority WBC} — classical Sentis \& Khatib formulation~\cite{Sentis2005}; no inequality constraint handling.
  \item \textbf{PD Joint Control} — decentralized joint-level PD; no whole-body coordination.
  \item \textbf{RL Policy} — pre-trained locomotion policy (e.g., Unitree H1 RL); used as agility upper-bound reference only.
\end{enumerate}

% ─── 5.4 SCENARIO S1 — STATIC BALANCING ─────────────────────
\section{S1: Static Balancing}
\label{sec:s1}

% Primary metric: CoM error, recovery time.
% Test: standing still under external pushes.

\lipsum[33]

\begin{table}[htb]
  \centering
  \caption{S1 — Static balancing results.}
  \label{tab:s1-results}
  \begin{tabular}{lccc}
    \toprule
    \textbf{Method} & $e_{\mathrm{CoM}}$ (cm) & $\eta_{\mathrm{FC}}$ (\%) & $t_{\mathrm{QP}}$ (ms) \\
    \midrule
    Proposed WBC     & --- & --- & --- \\
    Task-Priority WBC & --- & --- & --- \\
    PD Control       & --- & --- & --- \\
    \bottomrule
  \end{tabular}
\end{table}

\begin{criticalinsight}
  % TODO: Discuss if static balancing reveals any unexpected constraint
  % infeasibility or QP conditioning issues.
\end{criticalinsight}

% ─── 5.5 SCENARIO S2 — FLAT-GROUND WALKING ──────────────────
\section{S2: Flat-Ground Walking}
\label{sec:s2}

\lipsum[34]

\begin{criticalinsight}
  % TODO: Discuss CoM tracking performance during contact transitions.
\end{criticalinsight}

% ─── 5.6 SCENARIO S3 — STAIR CLIMBING ───────────────────────
\section{S3: Stair Climbing}
\label{sec:s3}

\lipsum[35]

\begin{criticalinsight}
  % TODO: Foot clearance and contact force analysis.
\end{criticalinsight}

% ─── 5.7 SCENARIO S4 — LOCO-MANIPULATION ───────────────────
\section{S4: Loco-Manipulation}
\label{sec:s4}

% Key scenario for H1 validation: simultaneous locomotion + manipulation.
% Metrics: end-effector error + CoM tracking.

\lipsum[36]

\begin{criticalinsight}
  % TODO: Quantify the coupling effect between locomotion and manipulation tasks.
  % Does task-priority assignment degrade either task? By how much?
\end{criticalinsight}

% ─── 5.8 SCENARIO S5 — PERTURBATION ROBUSTNESS ─────────────
\section{S5: Perturbation Robustness}
\label{sec:s5}

\lipsum[37]

\begin{criticalinsight}
  % TODO: Define the perturbation protocol precisely for reproducibility.
  % (Addresses §7.1 evaluation gap — perturbation protocols rarely reported.)
\end{criticalinsight}

% ─── 5.9 COMPARATIVE ANALYSIS ───────────────────────────────
\section{Comparative Analysis}
\label{sec:comparative}

% Synthesize results across all scenarios and baselines.
% Do NOT simply list numbers — argue what the data shows.

\lipsum[38]

\begin{criticalinsight}
  % TODO: Address anticipated reviewer criticism:
  % "RL baselines would outperform your WBC" — frame as expected (§18).
  % "Evaluation scenarios are too simple" — justify with loco-manipulation + perturbation.
\end{criticalinsight}
